\begin{lemma}[edge resolvent bounds, uniformly in the ensemble]\label{lem:edge}
Let $\eg\in(0,\frac1{10}]$, $N\ge N_1(\Par,\eg)$ \textup{(Lemma~\ref{lem:N0})}, $\eta=N^{-2/3-\eg}$, $\Psi=4N^{-1/3+2\eg}$, and
\[
I_{\rm edge}:=\big[\,2-N^{-2/3+\eg},\ 2+N^{-2/3+\eg}\,\big].
\]
Define the subset of $\Omega=\K^{\mathcal P_N}$
\begin{align*}
\Omega^c_{\rm edge}&:=\Big\{x\in\Omega:\ \big|G_{H(x)}(y+i\eta)_{kl}+\delta_{kl}\big|\le\tfrac12\Psi\ \ \text{for all }k,l\in\{1,\dots,N\}\text{ and all }y\in I_{\rm edge}\Big\},\\
\Omega_{\rm edge}&:=\Omega\setminus\Omega^c_{\rm edge}.
\end{align*}
Then $\Omega^c_{\rm edge}$ is closed, $\Omega_{\rm edge}$ is open (in particular both are Borel sets; they do not depend on the ensemble), and for every $\K\in\{\R,\C\}$ and
\emph{every} ensemble $M$ of size $N$ over $\K$ with $\fA(\Par)$:
\[
\P_M(\Omega_{\rm edge})\le N^{-20}.
\]
\end{lemma}
\begin{proof}
\emph{Step 0: topology.} Fix $k,l$ and $y\in I_{\rm edge}$, and put $z:=y+i\eta$. By Lemma~\ref{lem:basic}(h) and Lemma~\ref{lem:weyl}(d), for $x,x'\in\Omega$,
$|G_{H(x)}(z)_{kl}-G_{H(x')}(z)_{kl}|\le\eta^{-2}\|H(x)-H(x')\|\le\sqrt2\,\eta^{-2}|x-x'|$. So $x\mapsto G_{H(x)}(z)_{kl}$ is continuous and
$\{x:|G_{H(x)}(z)_{kl}+\delta_{kl}|\le\frac12\Psi\}$ is closed. $\Omega^c_{\rm edge}$ is the intersection of these sets over all $k,l$ and all $y\in I_{\rm edge}$, hence
closed; its complement $\Omega_{\rm edge}$ is open, hence Borel.

\emph{Step 1: the window lies in $S_{L_0}$.} Since $N\ge1$ and $-\frac23+\eg<0$, $N^{-2/3+\eg}\le1$, so $I_{\rm edge}\subset[1,3]\subset[-5,5]$. By \textup{(N1)},
$y+i\eta\in S_{L_0}$ for every $y\in I_{\rm edge}$, and \textup{(BB1)} holds with $L=L_0=A_0\log\log N$.

\emph{Step 2: the size of the control parameter in the window.} Let $y\in I_{\rm edge}$, $z=y+i\eta$. By Lemma~\ref{lem:mscB}(b),
\begin{equation}\label{eq:B3msc}
\im\msc(z)\le\sqrt2\,N^{-1/3+\eg/2},\qquad|\msc(z)+1|\le\sqrt2\,N^{-1/3+\eg/2}.
\end{equation}
Moreover $N\eta=N^{1/3-\eg}$. Hence
\[
\sqrt{\frac{\im\msc(z)}{N\eta}}\le\Big(\sqrt2\,N^{-\frac13+\frac{\eg}2}\cdot N^{-\frac13+\eg}\Big)^{1/2}=2^{1/4}N^{-\frac13+\frac34\eg}\le2N^{-\frac13+\eg},\qquad\frac1{N\eta}=N^{-\frac13+\eg},
\]
and therefore, by \textup{(N2)} (first inequality),
\begin{equation}\label{eq:B3ctrl}
(\log N)^{4L_0}\sqrt{\frac{\im\msc(z)}{N\eta}}+\frac{(\log N)^{4L_0}}{N\eta}\le3(\log N)^{4L_0}N^{-\frac13+\eg}\le N^{\eg/2}\cdot N^{-\frac13+\eg}=N^{-\frac13+\frac32\eg}.
\end{equation}

\emph{Step 3: $\Omega_{\rm edge}\subset U_{L_0}$.} Recall from \S\ref{sec:inputs} that $U_{L_0}$ is the union over $z\in S_{L_0}$ of the events
$\{\Lambda_d(z)+\Lambda_o(z)\ge(\log N)^{4L_0}\sqrt{\im\msc(z)/(N\eta)}+(\log N)^{4L_0}/(N\eta)\}$, where $\Lambda_d=\max_k|G_{kk}-\msc|$ and $\Lambda_o=\max_{k\ne l}|G_{kl}|$ \eqref{eq:defLambda},
$G=G_{H(x)}$. Let $x\in\Omega\setminus U_{L_0}$ and $y\in I_{\rm edge}$, $z=y+i\eta$. Since $z\in S_{L_0}$ (Step 1) and $x\notin U_{L_0}$, \eqref{eq:B3ctrl} gives
\[
\Lambda_d(z)+\Lambda_o(z)<N^{-\frac13+\frac32\eg}.
\]
Both $\Lambda_d$ and $\Lambda_o$ are nonnegative, so for all $k,l$: $|G_{kl}(z)-\delta_{kl}\msc(z)|\le\max(\Lambda_d(z),\Lambda_o(z))<N^{-1/3+3\eg/2}$. With \eqref{eq:B3msc},
\begin{align*}
\big|G_{kl}(z)+\delta_{kl}\big|&\le\big|G_{kl}(z)-\delta_{kl}\msc(z)\big|+\delta_{kl}\,|\msc(z)+1|\\
&<N^{-\frac13+\frac32\eg}+\sqrt2\,N^{-\frac13+\frac12\eg}=N^{-\frac13+2\eg}\Big(N^{-\frac{\eg}2}+\sqrt2\,N^{-\frac32\eg}\Big).
\end{align*}
By \textup{(N2)}, $N^{\eg/2}\ge3(\log N)^{4L_0}\ge3$ (recall $\log N\ge1$ and $L_0>0$), so $N^{-\eg/2}+\sqrt2N^{-3\eg/2}\le\frac13+\frac{\sqrt2}{27}<2$. Thus
$|G_{kl}(z)+\delta_{kl}|<2N^{-1/3+2\eg}=\frac12\Psi$ for all $k,l$ and all $y\in I_{\rm edge}$, i.e.\ $x\in\Omega^c_{\rm edge}$. We have shown
$\Omega\setminus U_{L_0}\subset\Omega^c_{\rm edge}$, that is, $\Omega_{\rm edge}\subset U_{L_0}$.

\emph{Step 4: probability.} $\Omega_{\rm edge}$ is a Borel subset of $U_{L_0}$ (Steps 0 and 3). Let $M$ be an ensemble of size $N$ with $\fA(\Par)$. By \textup{(BB1)} with $L=L_0$,
in the form stated in \S\ref{sec:inputs} (monotonicity of the outer probability), and by \textup{(N3)},
\[
\P_M(\Omega_{\rm edge})\le\P_M^*(U_{L_0})\le Ce^{-c(\log N)^{\phi L_0}}\le N^{-20}.\qedhere
\]
\end{proof}

\noindent\emph{Remarks.} (1) \emph{Exponent check.} The three contributions to $|G_{kl}+\delta_{kl}|$ are $(\log N)^{4L_0}(N\eta)^{-1}\approx N^{-1/3+\eg}$ (the dominant one),
$(\log N)^{4L_0}(\im\msc/(N\eta))^{1/2}\approx N^{-1/3+3\eg/4}$ and $|\msc+1|\approx N^{-1/3+\eg/2}$; the factor $(\log N)^{4L_0}\le\frac13N^{\eg/2}$ is absorbed by (N2), and the
remaining room $N^{\eg/2}$ up to $\frac12\Psi=2N^{-1/3+2\eg}$ is not used. The value $-1$ on the diagonal is the edge value $\msc(2)=-1$.

(2) \emph{Uniformity over the hybrids.} The constants $A_0,C,c,\phi,N_0$ of \textup{(BB1)}, hence $N_1(\Par,\eg)$, depend only on $\Par$ (and $\eg$), and every hybrid $M_k$,
$0\le k\le\Gamma$, satisfies $\fA(\Par)$ (Lemma~\ref{lem:hyb}(ii)). So the lemma holds for each $M_k$ with the same threshold and the same bound $N^{-20}$. The set $\Omega_{\rm edge}$
is the same for all ensembles; only its probability depends on $M$. Each Lindeberg step below uses the bound for one hybrid only; no union bound over $k$ is taken.

(3) \emph{Simultaneity in $y$.} The lemma controls all $y\in I_{\rm edge}$ at once; this is the simultaneity in $z\in S_L$ of Theorem~\ref{thm:21}(ii). By the remark~(4)
after Lemma~\ref{lem:expand}, the comparison branch would also work with the bound at two fixed points $y\in\{E_1,E_2\}$; we do not use this simplification.
(4) \emph{What is not used.} No bound for a matrix with an entry or a row removed is taken from Part~I: the matrix $Q=H^{[p]}$ of the next subsections does not satisfy (A),
and its resolvent is controlled deterministically by Lemma~\ref{lem:RfromS}(a).
